% =====================================================================
%  DATOS PERSONALES Y DE LA ACTIVIDAD  (rellenar antes de compilar)
% =====================================================================
\newcommand{\asignatura}{[Nombre de la asignatura]}
\newcommand{\actividad}{Diseño de un recurso educativo gamificado: \emph{Misión SQL}}
\newcommand{\bloque}{Bloque temático 3: Bases de datos}
\newcommand{\materiaFP}{Módulo Profesional 0484 \emph{Bases de datos} (1.\textsuperscript{er} curso de los CFGS de Desarrollo de Aplicaciones Multiplataforma y Desarrollo de Aplicaciones Web)}

\newcommand{\alumno}{[Nombre y apellidos]}
\newcommand{\dni}{[00000000X]}
\newcommand{\email}{[correo@alumno.uned.es]}
\newcommand{\centro}{[Centro asociado de referencia]}
\newcommand{\curso}{Curso 2026--2027}

% --- URL públicas de los recursos (sustituir cuando estén publicados) ---
\newcommand{\urlGenially}{https://view.genially.com/XXXXXXXX}
% Los tres juegos son páginas web propias (HTML + CSS + JavaScript)
% publicadas con GitHub Pages en un único repositorio.
\newcommand{\urlRepo}{https://github.com/USUARIO/mision-sql}
\newcommand{\urlJuegos}{https://USUARIO.github.io/mision-sql/}
\newcommand{\urlJuegoUno}{https://USUARIO.github.io/mision-sql/juego1/}
\newcommand{\urlJuegoDos}{https://USUARIO.github.io/mision-sql/juego2/}
\newcommand{\urlJuegoTres}{https://USUARIO.github.io/mision-sql/juego3/}

% --- Número de diapositivas de la presentación Genially ---
\newcommand{\numDiapositivas}{[XX]}
